\documentclass[journal]{IEEEtran}

\usepackage{iftex}
\ifTUTeX
  \usepackage{fontspec}
  \setmainfont{texgyretermes}[Extension=.otf, UprightFont=*-regular, BoldFont=*-bold,
    ItalicFont=*-italic, BoldItalicFont=*-bolditalic]
\else
  \usepackage[T1]{fontenc}
  \usepackage[utf8]{inputenc}
\fi
\usepackage{amsmath,amssymb}
\usepackage{booktabs}
\usepackage{tabularx}
\usepackage{array}
\usepackage{multirow}
\usepackage{xcolor}
\usepackage[hidelinks]{hyperref}
\usepackage{cite}

\newcolumntype{L}{>{\raggedright\arraybackslash}X}
\newcommand{\TWS}{\mathrm{TWS}}

\begin{document}

\title{Deep Learning and Diffusion Models for Enhancing GRACE/GRACE-FO Satellite Gravimetry: A Survey}

\author{H.~O\u{g}uz~Bolat%
\thanks{Draft of \today, prepared with Claude Code. Only works whose full text was obtained and read are cited; see Section~\ref{sec:method}.}}

\markboth{Draft}{Bolat: Deep Learning and Diffusion Models for GRACE Gravimetry}

\maketitle

\begin{abstract}
The GRACE and GRACE-FO missions have produced the only global, direct record of terrestrial water storage change since 2002. That record is coarse (an effective resolution of roughly 300\,km), noisy in structured ways, and interrupted by an eleven-month gap between the two missions. Machine learning is now used at almost every stage that follows the gravity solutions: filling gaps, reconstructing storage before 2002, downscaling to finer grids, forecasting, and correcting errors. In parallel, diffusion models have become the leading generative approach for downscaling, data assimilation and gap filling in weather, climate and ocean science. This survey brings the two literatures together. We review the measurement system and its error structure, the deep-learning methods applied to GRACE products and to the upstream processing chain, the foundations of diffusion models, and their applications in the geosciences. We restrict citations to works we have read in full. Three findings stand out. First, nearly all learned GRACE products are deterministic regressors whose fine-scale content is inherited from a hydrological-model or covariate prior; uncertainty, where reported, is per-pixel and uncalibrated. Second, none of the learned methods we read represents the actual GRACE observation operator or its correlated error covariance, and mass consistency is enforced only through soft penalties. Third, adjacent fields have already developed the needed pieces: residual diffusion, posterior sampling with explicit measurement operators, hard conservation constraints at inference time, and training from incomplete data. None of these has yet been applied to satellite gravimetry in the literature we could read. We close with a map from GRACE tasks to transferable diffusion techniques and a list of open problems.
\end{abstract}

\begin{IEEEkeywords}
GRACE, GRACE-FO, satellite gravimetry, terrestrial water storage, deep learning, diffusion models, downscaling, gap filling, uncertainty quantification.
\end{IEEEkeywords}

\section{Introduction}
\label{sec:intro}

\IEEEPARstart{S}{atellite} gravimetry from the Gravity Recovery and Climate Experiment (GRACE, 2002--2017) and its Follow-On (GRACE-FO, since 2018) measures month-to-month changes in Earth's gravity field. Over land, these changes are dominated by terrestrial water storage (TWS), the sum of groundwater, soil moisture, snow, surface water and glaciers. The record has shown regional freshwater trends worldwide~\cite{rodell2018emerging}, closed the sea-level budget together with altimetry and Argo~\cite{loomis2019regularization}, and traced ice-sheet mass loss across the change of mission~\cite{velicogna2020continuity}. It is now the only source of TWS as an Essential Climate Variable~\cite{rodell2023water}. Reviews of the mission's science and processing are given in~\cite{chen2022applications,rodell2023water}.

Three properties of the data limit how they can be used. The fields are coarse. The effective resolution is about 300\,km, and products delivered on 0.25$^\circ$ or 0.5$^\circ$ grids carry far less independent information than their pixel counts suggest~\cite{vishwakarma2018resolution,loomis2019regularization,bonin2026ocean}. The errors are structured. Spherical-harmonic solutions show north--south stripes from correlated errors, and the filters and regularisation that suppress the stripes also bias and leak the signal~\cite{kusche2007approximate,wiese2016leakage,loomis2019regularization}. And the record has holes: an eleven-month gap between the missions and scattered missing months before it~\cite{chen2022applications,landerer2020extending}.

Machine learning has been applied to all three. Convolutional networks learn the mismatch between land-surface models and GRACE~\cite{sun2019combining}. Bayesian networks bridge the inter-mission gap with uncertainty~\cite{mo2022bayesian}. Statistical and recurrent models reconstruct storage back to the early twentieth century~\cite{humphrey2019gracerec,yin2023gtwsmlrec,palazzoli2025graice}. Self-supervised encoder-decoders downscale GRACE to 0.5$^\circ$ globally~\cite{gou2024global}. Over the same period, diffusion models~\cite{ho2020ddpm,song2021scoresde} have become the dominant generative approach for downscaling, ensemble forecasting and data assimilation in atmospheric and ocean science~\cite{mardani2025corrdiff,price2025gencast,manshausen2025genda}. Their main appeal there is not lower point error but samples from a distribution: ensembles of spatially coherent fields whose spread can be calibrated.

This survey connects these literatures. Its questions are:
\begin{enumerate}
  \item What limits the information content of GRACE products, and what error models exist that a probabilistic method should respect?
  \item Which deep-learning methods have been applied to GRACE, at which stage of the processing chain, and how are they validated?
  \item Which diffusion-model techniques, developed in machine learning and in other geosciences, address the constraints that GRACE imposes: few training samples, no fine-scale truth, a non-trivial measurement operator, and a conservation requirement?
\end{enumerate}

\subsection{Review method}
\label{sec:method}

The literature was collected on 13 November 2026 by five parallel searches, one per sub-topic: GRACE measurement and processing; deep learning on GRACE products; machine learning in the gravimetry processing chain and complementary geodetic sensors; diffusion-model foundations; and diffusion and generative models in the Earth sciences. Sources were general web search, OpenAlex, Crossref, the arXiv API, the GFZ list of GRACE publications, Copernicus preprint servers and EGU abstracts. Each bibliographic record was checked against Crossref, OpenAlex or arXiv.

We apply one rule throughout: \emph{a work is cited only if its full text was obtained and read}. Text was extracted from each PDF, and every claim attributed to a paper here was checked against that text. Works known only from an abstract or a bibliographic record are not cited, however relevant they appear. This excludes a number of closed-access papers, and also open-access papers whose publishers blocked automated retrieval. Those are listed separately for retrieval and will be added in a later revision. The rule has consequences that the reader should keep in mind. Several canonical references for GRACE processing are absent here and are cited indirectly through the reviews that discuss them. The one diffusion-model study on GRACE products that our searches identified could only be read as an abstract and is therefore not discussed. Statements such as ``none of the methods reviewed'' refer to the works we read.

\subsection{Organisation}

Section~\ref{sec:grace} summarises the GRACE measurement system, its products, its effective resolution and its error models. Section~\ref{sec:mlgrace} reviews deep learning on gridded GRACE products by task. Section~\ref{sec:geodetic} covers machine learning upstream in the processing chain, complementary geodetic sensors and data assimilation. Section~\ref{sec:diffusion} introduces diffusion models with emphasis on conditioning and inverse problems. Section~\ref{sec:geodiff} reviews diffusion and other generative models in the Earth sciences. Section~\ref{sec:synthesis} maps GRACE tasks to transferable techniques and lists open problems. Section~\ref{sec:conclusion} concludes.

\section{GRACE and GRACE-FO: Measurement, Products and Limits}
\label{sec:grace}

\subsection{Measurement system}

GRACE and GRACE-FO fly two satellites in a near-polar orbit about 220\,km apart. A microwave ranging system measures changes in their separation. Each satellite carries an electrostatic accelerometer for non-gravitational forces, GPS receivers for precise orbits and star cameras for attitude~\cite{chen2022applications}. GRACE-FO adds a laser ranging interferometer (LRI). In orbit it reached a ranging noise of about 1\,nm\,Hz$^{-1/2}$ above 100\,mHz, far below the microwave system~\cite{abich2019inorbit}. Better ranging alone does not, however, yield much sharper gravity fields. A five-year closed-loop simulation found that the LRI improves unfiltered monthly solutions by about 23\% globally, but only by a few percent after decorrelation filtering. Accelerometer noise and errors in ocean-tide and non-tidal background models dominate the error budget~\cite{flechtner2016lri}. High-frequency atmosphere and ocean mass is removed with a de-aliasing product before inversion. Errors in that product propagate into the monthly fields as temporal aliasing. Adding a stochastic model of those errors to the inversion reduces noise by up to 25\%~\cite{kvas2019background}. The residual errors in Level-2 fields are therefore shaped by the instruments, the orbit and the background models; they are not white noise.

\subsection{Spherical harmonics and mascons}

Analysis centres invert the Level-1B data into monthly Stokes coefficients, typically to degree and order 60--96. Combination services merge solutions or normal equations from several centres and obtain lower noise than any single input~\cite{jaggi2020costg}. Spherical-harmonic (SH) solutions show north--south stripes, caused by correlated errors that follow from the along-track measurement geometry and from undersampled sub-monthly signals~\cite{chen2022applications}. They are suppressed by isotropic Gaussian smoothing, by fitting smooth polynomials to even and odd coefficients of each order, or by decorrelation filters. The DDK-type filters build a regularisation from a synthetic error covariance and a prior signal covariance. They remove stripes more efficiently than Gaussian filters while preserving amplitude. Kusche notes that isotropic smoothing ignores both latitude-dependent data density and anisotropic error correlations~\cite{kusche2007approximate}. Every filter attenuates the signal and leaks it across boundaries such as coastlines and basin divides. Products then apply restoration corrections~\cite{chen2022applications}.

Mascon solutions estimate surface mass in discrete parcels directly. They use prior information in the inversion to suppress correlated errors instead of filtering afterwards~\cite{wiese2016leakage,luthcke2013iterated,loomis2019regularization}. The JPL solution still needs a coastline filter, which reduces land--ocean leakage by about half in simulation, and gain factors for continental hydrology~\cite{wiese2016leakage}. The CSR solution uses about 41{,}000 hexagonal mascons of roughly 120\,km and is released on a 0.25$^\circ$ grid~\cite{bonin2026ocean}. The GSFC solution is estimated from range-rate data with full noise covariance~\cite{luthcke2013iterated}. Its later version derives the regularisation iteratively from data. Loomis et al.\ show that regularised mascons are biased by construction. Their total error must therefore combine a noise term with a leakage term computed from the monthly resolution operator~\cite{loomis2019regularization}. Regularisation acts as a prior. The resolution operator, which maps true mass to estimated mass, is the observation operator of the product.

\subsection{Effective resolution}

The delivery grid is not the resolution. CSR mascons are released at 0.25$^\circ$, but their true resolution is about 200--300\,km~\cite{bonin2026ocean}. The GSFC analysis recommends aggregating to about 520\,km for regional analysis outside Antarctica and about 270\,km over the Antarctic ice sheet~\cite{loomis2019regularization}. Rodell et al.\ give the resolvable area of monthly TWS as about 200{,}000\,km$^2$ at low latitudes and 90{,}000\,km$^2$ near the poles~\cite{rodell2018emerging}. A closed-loop study over 255 catchments found that, with the best correction schemes applied to SH products, catchments of about 63{,}000\,km$^2$ can be resolved at 2\,cm error. Filtered products with additive corrections need about 152{,}000\,km$^2$, and a 1\,cm error needs at least about 250{,}000\,km$^2$~\cite{vishwakarma2018resolution}. Any structure below roughly 300\,km in a ``high-resolution'' GRACE product must therefore come from somewhere other than the gravity measurement.

\subsection{Low degrees, geocentre and glacial isostatic adjustment}

GRACE determines the oblateness term $C_{20}$ poorly and cannot observe degree~1 (geocentre motion). Standard practice replaces $C_{20}$ with satellite laser ranging values and estimates degree~1 from GRACE combined with an ocean model~\cite{swenson2008geocenter,chen2022applications}. These conventions matter. A revised SLR-based $C_{20}$ changes Antarctic trends by $-15.4$\,Gt/yr and Greenland trends by $-3.5$\,Gt/yr relative to the earlier product~\cite{loomis2019oblateness}. Glacial isostatic adjustment must be removed with a model and is the dominant error term for Antarctic mass balance~\cite{chen2022applications}. A learned product inherits all of these choices from its training target.

\subsection{Gaps and the GRACE-FO transition}

The record has an eleven-month gap from July 2017 to May 2018~\cite{chen2022applications}. It also has scattered missing months late in the GRACE mission, caused by battery management. On GRACE-FO, the accelerometer of the second satellite showed correlated bias jumps and noise from 21 June 2018. Its data have since been replaced by a ``transplant'' from the first satellite with thruster modelling, the approach also used for the last seven GRACE months~\cite{landerer2020extending}. A physics-based alternative recovers the missing accelerometer signal from non-gravitational force models~\cite{behzadpour2021acc_recovery}. GRACE-FO nevertheless continues the GRACE record at equivalent precision~\cite{landerer2020extending}, and ice-sheet mass loss is continuous across the gap~\cite{velicogna2020continuity}. For a learned model, the gap months and the transplant months mark changes in data quality and error structure. Validation should be stratified across them.

\subsection{Error characterisation}

A probabilistic method should be judged against the error models that exist. Wahr et al.\ derived latitude- and smoothing-dependent error estimates for monthly mass~\cite{wahr2006accuracy}. Boergens et al.\ give an anisotropic, latitude-dependent covariance model for gridded TWS~\cite{boergens2020covariance}. Loomis et al.\ separate mascon noise from leakage~\cite{loomis2019regularization}. Triple collocation estimates GRACE land errors of about 11--12\,mm without a reference~\cite{vandijk2014reanalysis}. Over the ocean, GRACE bottom pressure agrees with in-situ recorders to about 1\,cm at low and middle latitudes and 1.5--2\,cm at high latitudes~\cite{chambers2012obp}. For closed-loop tests, the ESA Earth System Model provides 6-hourly mass fields with known truth~\cite{dobslaw2015esa_esm}, and a matching model of de-aliasing errors provides realistic perturbations~\cite{dobslaw2016dealiasing_errors}.

\subsection{Future missions}

GRACE-C is planned as a continuity mission in a polar orbit with laser ranging only. ESA's Next Generation Gravity Mission would add an inclined pair. Together they would form a Bender-type double-pair constellation (MAGIC), which doubles space-time sampling and adds east--west sensitivity against aliasing~\cite{haagmans2020nggm}. In a closed-loop study, the standard deviation of basin-average TWS noise in the Danube fell from 29.3\,mm for a simulated GRACE-C-like pair to 17.7\,mm for the inclined pair and 13.0\,mm for the constellation (preprint)~\cite{pfeffer2026nggm}. Cold-atom accelerometers coupled by a laser link improve simulated range-rate performance by 10--25\% between 1000 and 222\,km resolution~\cite{leveque2021quantum}. Methods trained on GRACE-like inputs will need a path to these richer inputs.

\section{Deep Learning on GRACE Products}
\label{sec:mlgrace}

Learned methods applied to gridded GRACE TWS anomalies fall into five task families: gap filling, reconstruction beyond the mission record, spatial and temporal downscaling, forecasting, and learned error correction. A sixth concern, uncertainty quantification, cuts across them. Nearly all of these methods share their predictors. Meteorological forcing (precipitation, temperature, evapotranspiration from ERA5, ERA5-Land or GLDAS) is almost universal. Storage simulated by a land-surface or global hydrological model is common, as are static fields such as topography and land cover~\cite{sun2019combining,mo2022bayesian,gou2024global}. Table~\ref{tab:mlgrace} lists representative studies.

\begin{table*}[!t]
\centering
\caption{Representative learned methods for GRACE TWS products (full texts read). UQ: form of uncertainty reported. Split: how evaluation months were chosen.}
\label{tab:mlgrace}
\footnotesize
\begin{tabularx}{\textwidth}{@{}l l L L l l@{}}
\toprule
Study & Task & Method & Prior / inputs & UQ & Split \\
\midrule
Sun et al.~\cite{sun2019combining} & model correction, gap & CNNs learn NOAH$-$GRACE mismatch & GLDAS-NOAH, P, T & none & temporal \\
Mo et al.~\cite{mo2022bayesian} & inter-mission gap & Bayesian CNN (SVGD, 20 particles) & ERA5-Land P, T, storage & per-pixel & temporal \\
Humphrey \& Gudmundsson~\cite{humphrey2019gracerec} & reconstruction 1901-- & statistical water-store model & P, T & 100-member ens. & detrended \\
Yin et al.~\cite{yin2023gtwsmlrec} & reconstruction 1940-- & RF (with MARS, GLM, ANN, GAM) & climate, land use, LAI & none & random \\
Palazzoli et al.~\cite{palazzoli2025graice} & reconstruction 1984-- & LSTM / BiLSTM & ERA5(-Land), SIF & none & temporal \\
Arzoumanidis et al.~\cite{arzoumanidis2026stgnn} & reconstruction 1940-- & spatio-temporal GNN & ERA5 P, ET, runoff & none & sequential \\
Uz et al.~\cite{uz2024highres} & hindcast 1994-- & residual conv.\ autoencoder & SLR fields, ERA5 & none & random \\
He et al.~\cite{he2021deep} & downscaling 11$\times$ & SRCNN, VDSR, RCAN & EALCO as truth & none & -- \\
Gou \& Soja~\cite{gou2024global} & downscaling to 0.5$^\circ$ & residual enc.-dec., dual loss & WGHM, GLDAS & deep ensemble & -- \\
Uz et al.~\cite{uz2026dynamic} & downscaling to $\sim$50\,km & variational U-Net, MC dropout & WGHM, ERA5 & per-pixel & -- \\
Li et al.~\cite{li2026physically} & downscaling 3$^\circ\!\to$0.1$^\circ$ & U-Net, conservation loss & PCR-GLOBWB 2 & seed ensemble & -- \\
Goracci et al.~\cite{goracci2026nggm} & downscaling to 1$^\circ$ & U-Net, aggregation loss only & ERA5-Land, topography & none & closed-loop \\
Dombos et al.~\cite{dombos2026dshift} & daily high-res fields & ESRGAN-type GAN & daily SH solutions & none & tuning years \\
Mo et al.~\cite{mo2025nearrealtime} & nowcast over latency & Bayesian CNN & past TWS, ERA5-Land & per-pixel & temporal \\
Steidl et al.~\cite{steidl2026hignn} & seasonal forecast & hierarchical GNN + LSTM & ERA5 & none & temporal \\
Chen et al.~\cite{chen2025efficient} & learned destriping & Bayesian-optimised CNN, attention & Level-2 SH & none & -- \\
\bottomrule
\end{tabularx}
\end{table*}

\subsection{Gap filling}

The first deep-learning study on GRACE did not predict GRACE directly. Sun et al.\ trained convolutional networks (VGG16, U-Net and SegNet variants) to predict the spatio-temporal mismatch between GLDAS-NOAH and GRACE over India. The corrected model output reached a country-average correlation of 0.94 and an NSE of 0.87. Because it no longer needs GRACE at prediction time, it also fills the inter-mission gap. The learned mismatch correlated with in-situ groundwater anomalies, so it largely represents the groundwater storage that the land-surface model lacks~\cite{sun2019combining}. Mo et al.\ filled the gap globally with a Bayesian CNN with residual connections and attention, trained by Stein variational gradient descent with an aleatoric noise term. It outperformed land-surface models and earlier gap products, most of all in arid regions, and returned predictive uncertainty~\cite{mo2022bayesian}.

Statistical gap fillers remain strong baselines. CIDR, an iterative singular-spectrum method with component-specific interpolation, improved relative RMS by 17.5--18.6\% and NSE by 27.6--42.6\% over standard SSA on artificial gaps in a 0.5$^\circ$ downscaled product~\cite{gao2025cidr}. A regional study split TWS into a trend, attributed to human activity, and a detrended part predicted by nine learners from model outputs. The explicit separation of trend from variability is its notable feature~\cite{mousavimehr2025nonstationary}.

\subsection{Reconstruction beyond the record}

Reconstructions extend TWS before 2002 using climate forcing. GRACE-REC fits a linear water-store model to detrended, deseasonalised GRACE and drives it with century-long precipitation and temperature. It provides 100-member ensembles for each of six forcing--product combinations~\cite{humphrey2019gracerec}. Trends and the seasonal cycle are removed by design. GTWS-MLrec compares five learners and reconstructs 1940--2022 at 0.25$^\circ$. It is evaluated against the land--ocean mass budget, 341 basin water budgets and over 10{,}000 streamflow gauges, but its train/test split is random in time~\cite{yin2023gtwsmlrec}. GRAiCE uses LSTM and BiLSTM networks with solar-induced fluorescence among its inputs and outperforms GRACE-REC against observed TWS~\cite{palazzoli2025graice}. A tree-based reconstruction selects predictors per grid cell with a Bayesian network. Extra-trees regression, not the CNN option, performed best in most cells~\cite{mandal2025reconstruction}. A spatio-temporal graph network with a fixed, interpretable adjacency reaches basin correlations within 0.025 of these products using only three forcing variables and a sequential split~\cite{arzoumanidis2026stgnn}. On the geodetic side, a residual convolutional autoencoder maps coarse SLR-derived mass fields and ERA5 variables to GRACE-like 1$^\circ$ grids from 1994~\cite{uz2024highres}. A conditional GAN calibrates WGHM against GRACE for long-term hindcasts, used as a deterministic generator~\cite{wang2024recgan}. Driven by climate alone, these reconstructions cannot represent human-driven trends such as groundwater depletion.

\subsection{Spatial and temporal downscaling}

\paragraph{Pixel-wise learning} The earliest downscaling methods fitted regressions (random forests, boosted trees, small networks) at GRACE scale and applied them to fine-scale predictors. Gains against independent data can be small. Over a crystalline aquifer in southern India, random-forest downscaling to 0.5$^\circ$ raised the mean correlation with 1006 wells from 0.76 to 0.79. The authors therefore propose measuring a ``downscaling gain'' against the non-downscaled product~\cite{pascal2022evaluating}. A partial-least-squares method conserves mass at catchment scale for 160 catchments. Its authors decline to rank input datasets because grid-scale validation is not possible~\cite{vishwakarma2021plsr}. A global 1\,km groundwater product chose gradient boosting over a deep network. It reaches $R^2 = 0.972$ when re-aggregated to GRACE scale but $R^2 = 0.538$ against held-out wells, and it shows a trend-sign mismatch in the High Plains aquifer for 2003--2013~\cite{zhang2026groundwater}.

\paragraph{Image super-resolution networks} He et al.\ transplanted SRCNN, VDSR and RCAN to downscale 110\,km JPL fields to 10\,km over Canada. A land-surface model served as high-resolution truth, and inputs were built by bicubic degradation. Re-aggregated RMSE was 14.4\,mm for VDSR, against 71.6\,mm for linear regression. The largest network (RCAN) underperformed the smaller VDSR, which the authors attribute to the 158 months of data~\cite{he2021deep}.

\paragraph{Self-supervised networks with a dual loss} The approach that has since become standard drops the model-as-truth assumption. Gou and Soja train a residual encoder-decoder with two losses. One requires agreement with GRACE after spatial averaging. The other requires correlation with WGHM fine-scale patterns. Their global 0.5$^\circ$ product has an area-weighted basin RMSE of 21.9\,mm against GRACE over 288 basins (49.2\,mm for WGHM), and deep ensembles give a median uncertainty of 7.3\,mm~\cite{gou2024global}. The same design downscales ocean bottom pressure to 0.25$^\circ$ by fusing two ocean reanalyses, partly to reduce dependence on a single prior~\cite{gou2025obp}. Uz et al.\ use a variational U-Net with Monte Carlo dropout and a dynamically weighted soft constraint between JPL mascons and WGHM/ERA5 patterns. They report a 95th-percentile uncertainty of about 3.2--3.3\,cm~\cite{uz2026dynamic}. Li et al.\ downscale from 3$^\circ$ to 0.1$^\circ$ over China with PCR-GLOBWB 2 as the pattern prior. They obtain a median correlation of 0.95 at 3$^\circ$ but 0.65 against wells where irrigation is intensive~\cite{li2026physically}. Goracci et al.\ keep only the aggregation-consistency loss, with fine structure supplied by ERA5-Land, topography and land cover. They state explicitly that the fine-scale redistribution remains under-constrained. Using simulated NGGM and MAGIC products, they show that the same network produces better 1$^\circ$ fields from richer inputs~\cite{goracci2026nggm}.

\paragraph{Adversarial and temporal downscaling} D-SHIFT trains an ESRGAN-type generator with a U-Net discriminator to transfer mascon-level spatial detail to daily SH solutions. Its monthly global RMSE is about 2.3\,cm. The authors report increased high-frequency noise and isolated daily excursions, there is no ensemble, and the held-out years were also used for tuning~\cite{dombos2026dshift}.

\paragraph{Independent evaluations} Intercomparisons by third parties are beginning. Over 288 basins, the median RMSE against GRACE is 1.87\,cm for the Gou--Soja product, 2.24\,cm for a non-learned joint inversion and 3.40\,cm for an assimilation product. The joint-inversion authors describe their own product as a GRACE-constrained reconstruction informed by high-resolution spatial priors~\cite{xiong2026sysu}. A well-based evaluation over India found that a mass-conserving statistical method improved on the mascon product in 17 of 22 mascons with wells, and the Gou--Soja product in 16~\cite{suryawanshi2026efficacy}. Pakrashi et al.\ take a statistical rather than learned route. A Bayesian hierarchical model links CSR mascon blocks and well observations to one latent fine-scale Gaussian Markov random field and returns posterior standard-deviation maps~\cite{pakrashi2026highplains}. It is the only study we read that writes an explicit areal observation model for GRACE and returns a posterior.

\subsection{Forecasting and nowcasting}

GRACE products appear with a latency of about three months. Mo et al.\ predict TWS over this latency from twelve past months and ERA5-Land forcing with a Bayesian CNN. They obtain median correlations of 0.92--0.95 one to three months ahead and detect droughts before GRACE data are released~\cite{mo2025nearrealtime}. Because forcing for the target months is an input, this is nowcasting rather than forecasting. Steidl et al.\ forecast up to six months with a hierarchical graph network and an LSTM trained on a reconstruction. One-month RMSE is 1.83\,cm against 3.70\,cm for climatology, and skill is lost by six months~\cite{steidl2026hignn}. Basin-scale studies compare backpropagation networks, LSTMs and BiLSTMs with attention, with no consistent winner across regions~\cite{lu2024changes}. A physics-constrained LSTM separates groundwater from GRACE TWS. It bounds trends by the moisture supply available from the atmosphere and enforces basin mass balance during training~\cite{sarkar2026physics}.

\subsection{Learned error correction and signal separation}

Learning directly on Level-2 data is rarer. Chen et al.\ train a Bayesian-optimised CNN with spatial attention to map Level-2 data to TWS, using processed Level-3 and mascon fields as targets. They reach a test NSE of 0.60 and suggest recovery from Level-1 data as future work~\cite{chen2025efficient}. Heller-Kaikov et al.\ compare a conditional GAN with principal component analysis for separating atmosphere--ocean from hydrological signals in closed-loop synthetic data. The two achieve comparable errors~\cite{hellerkaikov2023signal}.

\subsection{Uncertainty quantification}

Uncertainty, where reported, comes from Bayesian CNN particles~\cite{mo2022bayesian,mo2025nearrealtime}, deep ensembles~\cite{gou2024global}, Monte Carlo dropout~\cite{uz2026dynamic}, seed ensembles~\cite{li2026physically} or statistical ensembles~\cite{humphrey2019gracerec}. All of these are summarised as per-pixel spread. None of the deep-learning studies we read draws spatially coherent sample fields or tests calibration with proper scores, spread--skill ratios or rank histograms. A survey of uncertainty quantification for deep-learning climate variables, with GRACE as a worked example, shows why this matters. In its application cases, ignoring input uncertainty strongly underestimates total uncertainty, and ignoring model uncertainty underestimates it as well~\cite{gou2026uncertainty}.

\subsection{Validation practice and recurring critiques}

Five issues recur across the studies read.
\begin{enumerate}
  \item \emph{Consistency at GRACE scale is enforced, not tested.} Re-aggregation scores are high because the loss requires them~\cite{goracci2026nggm,zhang2026groundwater}. Independent data are needed below 300\,km, and there the scores fall~\cite{zhang2026groundwater,li2026physically}.
  \item \emph{Fine scale comes from the prior.} WGHM, PCR-GLOBWB or ERA5 supply the sub-GRACE structure~\cite{gou2024global,uz2026dynamic,li2026physically,xiong2026sysu}. Hydrological models are weakest where GRACE matters most: irrigation, groundwater abstraction, reservoirs and glaciers~\cite{springer2026review}.
  \item \emph{Mass consistency is soft.} Every network listed in Table~\ref{tab:mlgrace} that enforces agreement with GRACE does so through a loss penalty.
  \item \emph{The observation operator is a block average.} No learned method represents SH truncation, the filter or the mascon resolution operator, or the correlated error covariance.
  \item \emph{Small samples and leaky splits.} With about 250 monthly fields, random month splits~\cite{yin2023gtwsmlrec,uz2024highres} and reuse of tuning years for evaluation~\cite{dombos2026dshift} inflate reported skill. Sequential splits are used in~\cite{arzoumanidis2026stgnn,lu2024changes,steidl2026hignn}.
\end{enumerate}
A review of GRACE data assimilation names the resolution mismatch between GRACE and models as an open problem, and finds no consensus on the best assimilation approach~\cite{springer2026review}.

\section{Machine Learning in the Gravimetry Chain, Complementary Sensors and Assimilation}
\label{sec:geodetic}

The processing chain runs from Level-1 instrument data, through Level-2 harmonics, to post-processed grids. Machine learning has so far touched it mostly at its edges.

\subsection{Level-1 instrument data}

The degraded GRACE-FO accelerometer made accelerometer transplant an operational need. The operational transplant leaves $C_{20}$ and $C_{30}$ sensitive to how it is done. A physics-based recovery from non-gravitational force models reduces noise in the monthly fields and improves both coefficients~\cite{behzadpour2021acc_recovery}. Learned versions are emerging. Multilayer perceptrons, LSTMs and BiLSTMs transplant accelerometer data in closed-loop simulations of a future mission with cold-atom sensors~\cite{romeshkani2025ml_transplant}. Self-supervised Noise2Noise training denoises GRACE-FO and Taiji-1 accelerometer data~\cite{yang2023noise2noise_acc}. We found no learned treatment of ranging noise.

\subsection{Level-2 filtering and mascon regularisation}

The classical filters described in Section~\ref{sec:grace} and the hand-designed mascon regularisations of the analysis centres~\cite{wiese2016leakage,loomis2019regularization} remain standard. Learned work at Level-2 is limited to signal separation~\cite{hellerkaikov2023signal} and to the Level-2-to-TWS mapping of~\cite{chen2025efficient}. Loomis et al.\ show that a regularised mascon solution is biased and that its error is noise plus leakage~\cite{loomis2019regularization}. A learned prior would have to be characterised in the same terms. We found no learned prior used inside a mascon or SH inversion.

\subsection{Neural representations of the gravity field}

Physics-informed neural networks represent the gravitational potential as a continuous function constrained by Laplace's equation. The third generation of these models addresses feature divergence, low-altitude sample bias, numerical instability and extrapolation error~\cite{martin2024pinngm3}. Neural density fields learn the mass distribution of small bodies from their gravity signal~\cite{izzo2022geodesynets}. Both have been applied to static fields only.

\subsection{Low degrees and glacial isostatic adjustment}

Degree~1 is estimated from GRACE and an ocean model~\cite{swenson2008geocenter}, and $C_{20}$ from laser ranging~\cite{loomis2019oblateness}. No learned estimator appeared among the works we read. For glacial isostatic adjustment, neural networks have been trained as fast surrogates of forward Earth-deformation models~\cite{love2024gia_surrogate}, not as estimators from GRACE data.

\subsection{Complementary sensors}

\paragraph{Swarm} GPS-derived gravity from the Swarm satellites carries information to roughly spherical-harmonic degree~12. It needs about 750\,km smoothing to agree with GRACE at about 4\,cm over land~\cite{teixeira2020swarm_gfm,luck2018swarm_ocean}. Projecting Swarm fields onto GRACE empirical orthogonal functions raises their effective resolution~\cite{richter2021swarm_reconstruction}.

\paragraph{Satellite laser ranging} EOF-based hybrid SLR fields extend low-resolution mass change before 2002~\cite{locher2021slr_hybrid}. The residual autoencoder of~\cite{uz2024highres} learns the mapping from SLR to GRACE resolution.

\paragraph{GNSS loading} Vertical displacements of GNSS stations respond elastically to surface loading and give an independent constraint on TWS at finer scales than GRACE. Carlson et al.\ invert GNSS and GRACE jointly with a wavelet decomposition~\cite{carlson2022gnss_grace}. A 3D CNN that combines 41 GNSS stations with GRACE outperforms a load Green's-function inversion in reproducing the observed displacements~\cite{wei2025cnn_gnss_loading}.

\paragraph{Multi-sensor reanalysis} A global water-cycle reanalysis for 2003--2012 merges GRACE and water-level altimetry with a multi-model hydrological ensemble. It estimates errors by triple collocation: 11--12\,mm for GRACE storage over land against 17--28\,mm for model storage changes~\cite{vandijk2014reanalysis}.

\subsection{Data assimilation and hybrid models}

GRACE is assimilated into land-surface and hydrological models to downscale it and to partition storage. A recent review of about 60 such studies finds the resolution mismatch unresolved, along with the choice of assimilation scheme~\cite{springer2026review}. Learned components appear in three ways. Networks learn the residual between a model and GRACE~\cite{sun2019combining}. Self-supervised networks act as the assimilation step~\cite{gou2024global,gou2025obp}. And neural components are embedded in differentiable hydrological models trained against GRACE together with runoff, evapotranspiration and snow~\cite{kraft2022h2m,baghirov2025h2mv}. A review gives the general case for differentiable modelling~\cite{shen2023differentiable}. In all of these, GRACE acts as a weak, coarse-scale constraint on a model that supplies the fine scales.

\subsection{Benchmarks}

No community benchmark exists for learned gravimetry processing. The closest resource is closed-loop simulation. The ESA Earth System Model supplies 6-hourly mass fields with known truth~\cite{dobslaw2015esa_esm}. A de-aliasing error model supplies realistic perturbations~\cite{dobslaw2016dealiasing_errors}. Mission simulators produce the corresponding noisy solutions~\cite{pfeffer2026nggm,goracci2026nggm}. Real-data checks can use GNSS loading~\cite{carlson2022gnss_grace}, ocean-bottom pressure~\cite{chambers2012obp}, ice-sheet mass budgets~\cite{velicogna2020continuity} and wells~\cite{pascal2022evaluating,suryawanshi2026efficacy}. Calibration can be checked against formal covariance and error models~\cite{wahr2006accuracy,boergens2020covariance,loomis2019regularization}.

\section{Diffusion Models: Foundations Relevant to Gravimetry}
\label{sec:diffusion}

This section introduces the parts of the diffusion-model literature that bear on GRACE: the training objective, conditioning, posterior sampling with a known measurement operator, restoration-specific processes, spherical geometry and uncertainty. Broader surveys are~\cite{yang2023survey,croitoru2023survey}; inverse problems are surveyed in~\cite{daras2024survey}.

\subsection{Forward and reverse processes}

A diffusion model gradually corrupts data $x_0 \sim p_\mathrm{data}$ with noise and learns to reverse the corruption~\cite{sohldickstein2015deep}. In the discrete formulation of DDPM, $q(x_t \mid x_0) = \mathcal{N}(\sqrt{\bar\alpha_t}\,x_0, (1-\bar\alpha_t) I)$, and a network trained to predict the added noise defines the reverse chain~\cite{ho2020ddpm}. The continuous formulation writes the corruption as a stochastic differential equation $\mathrm{d}x = f(x,t)\,\mathrm{d}t + g(t)\,\mathrm{d}w$. Its reverse-time SDE and a deterministic probability-flow ODE with the same marginals both depend on the data only through the score $\nabla_x \log p_t(x)$~\cite{song2021scoresde}. The score is learned by denoising score matching, a regression of a network on noisy copies of the data~\cite{hyvarinen2005score,song2019ncsn}. Karras et al.\ separate the design choices of noise schedule, network preconditioning, loss weighting and sampler. Their EDM recipe, with a second-order Heun sampler of about 18--50 steps, is now the usual starting point in scientific applications~\cite{karras2022edm,mardani2025corrdiff,price2025gencast}. Sampling cost is reduced by implicit samplers~\cite{song2021ddim}, learned variances~\cite{nichol2021improved}, dedicated ODE solvers~\cite{lu2022dpmsolver}, progressive distillation~\cite{salimans2022progressive} and consistency models that sample in one or a few steps~\cite{song2023consistency}. Latent diffusion moves the process into the compressed space of an autoencoder~\cite{rombach2022ldm}.

Flow matching~\cite{lipman2023flow}, rectified flow~\cite{liu2023rectified} and stochastic interpolants~\cite{albergo2023interpolants} regress a velocity field along prescribed paths between two distributions; diffusion is a special case. These formulations allow a non-Gaussian source. That matters for restoration, where the natural source is the degraded observation itself.

\subsection{Conditioning}

There are three ways to make samples depend on an observation $y$.

\paragraph{Conditional training} The network receives $y$ as an extra input, by channel concatenation or cross-attention, and is trained on pairs $(x, y)$. This is the approach of conditional score models~\cite{batzolis2021conditional}, image super-resolution (SR3)~\cite{saharia2023sr3}, general image-to-image translation (Palette)~\cite{saharia2022palette} and probabilistic time-series imputation (CSDI)~\cite{tashiro2021csdi}.

\paragraph{Guidance} Classifier guidance adds the gradient of a classifier's log-likelihood to the score~\cite{dhariwal2021beat}. Classifier-free guidance mixes conditional and unconditional scores from one network~\cite{ho2022cfg}.

\paragraph{Training-free posterior sampling} An unconditional prior is combined with a likelihood at inference time using Bayes' rule, $\nabla \log p_t(x \mid y) = \nabla \log p_t(x) + \nabla \log p_t(y \mid x)$. The second term is intractable for $t>0$, and methods differ in how they approximate it~\cite{daras2024survey}. Replacement and projection methods overwrite the observed components at each step. RePaint does this for inpainting~\cite{lugmayr2022repaint}, and DDNM decomposes the solution into the range and null space of a linear operator~\cite{wang2023ddnm}. SNIPS and DDRM handle a noisy linear operator $y = Ax + n$ in the basis of its singular value decomposition~\cite{kawar2021snips,kawar2022ddrm}. Manifold-constrained gradients and diffusion posterior sampling (DPS) use the gradient of the likelihood evaluated at the denoised estimate. They apply to noisy and even nonlinear differentiable operators~\cite{chung2022mcg,chung2023dps}. Pseudoinverse guidance ($\Pi$GDM) uses a Gaussian approximation of that likelihood~\cite{song2023pigdm}. Plug-and-play splitting~\cite{zhu2023diffpir} and variational formulations~\cite{mardani2024reddiff} are alternatives. Score-based models were applied early to sparse-view tomography in medical imaging~\cite{song2022medical}. SDEdit is a simpler relative. It noises an observation to an intermediate level and denoises it with the prior~\cite{meng2022sdedit}.

Posterior sampling has a direct reading for GRACE. A monthly solution is, to first order, a linear functional of the true mass field: truncation to a maximum degree, a filter or a mascon averaging kernel, plus correlated noise whose covariance is partly known (Section~\ref{sec:grace}). A prior trained on realistic storage fields could therefore be combined with the operator and covariance of a specific product at inference, without retraining for each filter or release. Because SH truncation is diagonal in the harmonic basis, the SVD formulation of DDRM is particularly natural. These samplers approximate the posterior, so their ensembles are not guaranteed to be calibrated, and calibration must be checked~\cite{daras2024survey}.

\subsection{Processes built for restoration}

Several methods replace pure noise with the degradation itself. Schrödinger bridges connect two arbitrary distributions~\cite{debortoli2021dsb}. I$^2$SB builds the bridge between degraded and clean images and starts sampling from the degraded input~\cite{liu2023i2sb}. Mean-reverting SDEs~\cite{luo2023irsde}, ResShift's residual shifting~\cite{yue2023resshift} and direct iteration (InDI)~\cite{delbracio2023indi} follow the same idea. Cold Diffusion inverts deterministic degradations such as blurring~\cite{bansal2023cold}, and inverse heat dissipation uses blurring as the generative process~\cite{rissanen2023heat}. Smoothing and truncation in the spectral domain are close relatives of the Gaussian and decorrelation filters applied to GRACE.

\subsection{Spatio-temporal models}

Video diffusion factorises attention over space and time~\cite{ho2022video}. DYffusion couples a learned temporal interpolator with the diffusion process for forecasting~\cite{ruhlingcachay2023dyffusion}. PDE-Refiner shows that solvers trained with mean-squared error neglect low-amplitude, often high-frequency, components, which limits rollout accuracy. Diffusion-style refinement recovers them~\cite{lippe2023pderefiner}. The same argument applies to the spectral power that GRACE filtering removes. Score-based data assimilation trains a score over short trajectory segments and conditions on arbitrary observations at inference~\cite{rozet2023sda}.

\subsection{Spherical geometry}

GRACE fields live on the sphere. Equirectangular grids with standard U-Nets are simple but distort the poles. Spherical CNNs are equivariant to rotation~\cite{cohen2018spherical}. Spherical Fourier neural operators convolve in the harmonic domain, which suits SH-truncated data~\cite{bonev2023sfno}. Icosahedral graph meshes are used by GenCast~\cite{price2025gencast}. Riemannian diffusion models generate points on manifolds~\cite{debortoli2022riemannian,huang2022riemannian}. Gridded TWS, however, is a field over the sphere, which manifold diffusion fields address~\cite{elhag2024mdf}.

\subsection{Uncertainty and evaluation}

Diffusion samples form an ensemble, but the ensemble's spread need not match the error. Weather applications evaluate calibration with the continuous ranked probability score, spread--skill ratios and rank histograms~\cite{price2025gencast}. Conffusion~\cite{horwitz2022conffusion} and BayesDiff~\cite{kou2024bayesdiff} estimate pixel-wise intervals or variances. Geoscience applications show these ensembles are often under-dispersive (Section~\ref{sec:geodiff}).

\section{Diffusion and Generative Models in the Earth Sciences}
\label{sec:geodiff}

Generative models for gridded geophysical fields cluster into six families: downscaling, physically constrained generation, probabilistic forecasting, data assimilation and reconstruction from sparse observations, gap filling, and climate emulation. A review of diffusion in remote sensing is~\cite{liu2024diffrs}; one for time-series and spatio-temporal data is~\cite{yang2024diffsurvey}. Table~\ref{tab:geodiff} lists the designs most relevant to GRACE.

\begin{table*}[!t]
\centering
\caption{Diffusion-model designs in the Earth sciences with direct analogues for GRACE (full texts read).}
\label{tab:geodiff}
\footnotesize
\begin{tabularx}{\textwidth}{@{}l L L L@{}}
\toprule
Study & Design & Application & GRACE analogue \\
\midrule
CorrDiff~\cite{mardani2025corrdiff} & regression mean + EDM diffusion on residual & 25\,km$\to$2\,km weather, Taiwan & residual diffusion on a deterministic downscaler \\
R2-D2~\cite{lopezgomez2025r2d2} & regional model, then diffusion on residual & CMIP6 ensembles to 9\,km & GRACE-constrained hydrology + generative refinement \\
Tie et al.~\cite{tie2026zssd} & unconditional prior + DPS with area-weighted averaging operator & zero-shot GCM downscaling & guidance through the GRACE averaging kernel \\
Hess et al.~\cite{hess2025scaleadaptive} & consistency model; noise level at PSD intersection & ESM precipitation & trust scales above $\sim$300\,km, generate below \\
Aich et al.~\cite{aich2026conddiff} & shared embedding by smoothing above a correction scale & bias correction + downscaling & preserve GRACE-resolved scales \\
Bochow et al.~\cite{bochow2026smb} & consistency model + hard conservation at inference & Greenland SMB, 32$\times$ & mass-conserving monthly downscaling \\
Forestier et al.~\cite{forestier2026sss} & prior + pseudo-inverse guidance + high-res auxiliaries & sea-surface salinity & GRACE + covariates \\
Manshausen et al.~\cite{manshausen2025genda} & unconditional EDM + score-based DA of stations & km-scale weather & GRACE + wells / GNSS fusion \\
GenCRest~\cite{qian2026gencrest} & prior on monthly fields + likelihood guidance with obs.\ error & historical climate reconstruction & pre-2002 reconstruction, gap \\
Barth et al.~\cite{barth2024dindiff} & DDPM trained on incomplete images & chlorophyll gap filling & no gap-free fine-scale truth \\
DiffESM~\cite{bassetti2024diffesm} & space--time diffusion conditioned on monthly means & ESM emulation & monthly-to-daily disaggregation \\
\bottomrule
\end{tabularx}
\end{table*}

\subsection{Downscaling}

\paragraph{Residual diffusion} CorrDiff introduced the design now dominant in weather downscaling. A U-Net predicts the conditional mean, and an EDM diffusion model samples the residual. Directly conditioned diffusion had converged slowly and produced incoherent structure at a 12.5$\times$ ratio. Trained on three years of hourly data, CorrDiff recovers kinetic-energy spectra at 10--200\,km and radar-reflectivity distributions. Its temperature tails improve only incrementally, and its ensembles are under-dispersive for most variables~\cite{mardani2025corrdiff}. R2-D2 replaces the regression stage with a regional climate model and samples the residual between 9\,km fields and interpolated 45\,km fields. It lowers CRPS by more than 40\% relative to a standard statistical method, with spread--skill ratios of 0.66--0.80~\cite{lopezgomez2025r2d2}. A residual latent diffusion model over Italy beats a GAN with the same backbone and reports its cost: about 870--1100 GPU-hours of training against 250--380 for the U-Net~\cite{tomasi2025ldm}. Video diffusion on residuals yields temporally coherent fine scales~\cite{srivastava2024stvd}. ArchesWeatherGen identifies the main obstacle of residual designs: the deterministic model overfits. The remedy is to train the residual model on data the deterministic model has not seen~\cite{couairon2024arches}. With about 250 GRACE months this will be the binding constraint.

\paragraph{Directly conditioned models} Diffusion downscalers outperform U-Nets with the same architecture~\cite{watt2024diffdownscaling}. One reconstructs 180 years of monthly climate over East Asia, although it was trained on 6-hourly data and applied to monthly means~\cite{ling2024dpdm}. Precipitation emulators produce realistic extremes up to 100-year return periods and reach good climatology with 14 years of daily training data~\cite{addison2026cpmgem}. Conditioning design matters. Channel concatenation gives the best pointwise CRPS but over-smooths extremes, while cross-attention to a frozen foundation-model encoder matches full-data skill with five years of data~\cite{ribeiro2026prithvi}. Heterogeneous satellite and station inputs can be fused by cross-attention and used as guidance at sampling time~\cite{tu2025mods}. A physically guided model for land-surface temperature produces ensemble spread that correlates linearly with error, and it examines which aggregation operator links fine to coarse fields~\cite{zhang2025pgdm}.

\paragraph{Priors with guidance} When paired data do not exist, the prior is trained on fine-scale fields alone and the coarse observation is imposed at sampling. Bischoff and Deck noise the coarse field until small-scale information is lost, then denoise with a separately trained high-resolution model~\cite{bischoff2024bridges}. Hess et al.\ choose that noise level where the power spectra of the coarse and fine data intersect, and denoise in one step with a consistency model~\cite{hess2025scaleadaptive}. Aich et al.\ map observations and model fields to a shared space by smoothing above a correction scale, which preserves the large scales~\cite{aich2026conddiff}. Wan et al.\ first remove large-scale bias by optimal transport~\cite{wan2023debias}. With an explicit operator, Tie et al.\ guide a reanalysis-trained prior through an area-weighted average. Supervised baselines trained on reanalysis pairs, including CorrDiff, degrade on real climate-model inputs to errors worse than the coarse input; the zero-shot posterior sampler generalises~\cite{tie2026zssd}. Coarse grids and sparse radiosondes have been fused in one Bayesian sampler~\cite{chakraborty2025multimodal}. For sea-surface salinity, a multichannel prior guided toward the coarse field and high-resolution auxiliaries is comparable in accuracy to a strong CNN, and adds an ensemble whose calibration is not established~\cite{forestier2026sss}.

\paragraph{Conditional versus guided} Suresh Babu et al.\ compare the two families directly on turbulence. Conditionally trained models reconstruct fine scales and tails, and their spread predicts error. DPS gives smoothed fields and SDEdit gives unphysical ones, and neither guided method carries information into unobserved regions~\cite{sureshbabu2025guided}. Read together with~\cite{tie2026zssd}, this suggests a trade-off: conditional models are more accurate in-distribution, while guided priors are more robust to domain shift. GRACE faces both a shortage of real pairs and a shift between model-based training data and observations.

\subsection{Physical constraints and conservation}

Harder et al.\ design output layers that reproduce every coarse-cell mean exactly. Such hard constraints usually improve accuracy. Soft penalties leave violations and can raise error, for example RMSE from 0.657 to 0.801 in one configuration. The exception is when coarse and fine data come from inconsistent sources~\cite{harder2023hard}. A constrained GAN that conserves the model's global mean cannot improve that mean against observations, so the constraint must be tied to the observed quantity~\cite{hess2022pcgan}. Constraints can also be enforced during sampling without retraining. Bochow et al.\ apply a pooling constraint at inference when downscaling Greenland surface mass balance by 32$\times$~\cite{bochow2026smb}. Projected diffusion certifies convex constraints by repeated projection~\cite{christopher2024projected}. ECI sampling enforces hard constraints exactly and without gradients in flow-matching models~\cite{cheng2025eci}. For a linear averaging operator, projection is exact and cheap. Physics can instead enter as a training loss on PDE residuals~\cite{bastek2025pidm}, as residual-gradient steps during sampling~\cite{jacobsen2023cocogen,shu2023pidiffusion}, or as joint generation of states and coefficients under partial observation~\cite{huang2024diffusionpde}. A diffusion model trained only on high-fidelity flow data is much more robust than a direct-mapping CNN to changes in the low-fidelity input~\cite{shu2023pidiffusion}. GRACE inputs change between analysis centres and releases. Some models conserve without enforcement. ACE nearly conserves column moisture but has a surface-pressure bias of about $-25$\,Pa~\cite{wattmeyer2023ace}.

\subsection{Probabilistic forecasting and ensembles}

GenCast produces 15-day global ensembles at 0.25$^\circ$ with a graph-transformer on an icosahedral mesh. It beats the operational ECMWF ensemble on 97.4\% of 1320 targets in CRPS~\cite{price2025gencast}. SEEDS amplifies one or two physics-based members into hundreds. Its post-processing mode, trained toward reanalysis, corrects bias and improves extreme-event probabilities~\cite{li2024seeds}. That is a template for turning a handful of GRACE solutions from different centres or filters into a calibrated ensemble. StormCast emulates a convection-allowing model and remains robust when its conditioning is swapped at test time~\cite{pathak2026stormcast}. A VAE-based system provides a non-diffusion probabilistic baseline~\cite{zhong2025fuxiens}. In nowcasting, a GAN fixed the blurriness of deterministic predictions~\cite{ravuri2021dgmr}. Latent diffusion then improved uncertainty calibration over it~\cite{leinonen2023ldcast}, and cascaded deterministic--generative designs improved extremes~\cite{gong2024cascast}. Constraint-aware latent diffusion~\cite{gao2023prediff} adjusts transitions toward soft constraints. Data-space and latent diffusion of Arctic sea ice reach similar scores, but the latent version produces smoother fields and loses small-scale information~\cite{finn2024seaice}. GRACE grids are small, so data-space diffusion is affordable.

\subsection{Data assimilation and reconstruction from sparse observations}

Score-based data assimilation conditions a trajectory-level prior on observations at inference, so a new observation operator needs no retraining~\cite{rozet2023sda}. Scaled to a two-layer quasi-geostrophic model, it recovers trajectories from sparse stations but with less small-scale detail than the simulations~\cite{rozet2023sdaqg}. DiffDA repurposes a forecast network as a denoiser for global analyses~\cite{huang2024diffda}. Appa handles reanalysis, filtering and forecasting in one latent model, but reports degraded fine scales and unvalidated posterior calibration~\cite{andry2025appa}. Multimodal latent assimilation is more robust than pixel-space assimilation~\cite{qu2024deepgenda}. A km-scale weather prior steered by a few dozen stations beats the operational analysis on held-out stations and reconstructs unobserved variables~\cite{manshausen2025genda}. In the ocean, Lagrangian observations at 99.9\% sparsity are inverted by a diffusion model conditioned on a neural-operator prediction~\cite{asefi2025lagrangian}. Global sea-surface temperature is reconstructed by fusing satellite and in-situ data progressively~\cite{li2026sage}, and subsurface temperature from sparse profiles with up to 92.5\% missing data~\cite{song2025reconmost}. GenCRest is the closest analogue to pre-2002 TWS reconstruction. It is pretrained on monthly fields from climate models, then trained on reanalysis. It reconstructs historical temperature and precipitation with likelihood guidance that includes observational error covariances, and returns ensembles of plausible histories~\cite{qian2026gencrest}.

\subsection{Gap filling}

Barth et al.\ train a denoising diffusion model on incomplete satellite chlorophyll images without any gap-free truth. It beats a convolutional baseline, preserves small scales and gives a reliable ensemble spread~\cite{barth2024dindiff}. DIRECT applies rectified flow to sparse sea-surface temperature and reduces error in occluded regions by 6--14\% relative to the CRITER reconstruction, with calibrated uncertainty~\cite{rovscek2026direct}. Precipitation maps are inpainted as video~\cite{kishikawa2025precipinpaint}. In geodesy and geophysics, a diffusion model interpolates and denoises exploration gravity and magnetic grids. It is the only gravity-field diffusion study we found, but it concerns local surveys~\cite{hou2025gravmag}. Another conditions bathymetry gap filling on altimetry-derived gravity anomalies~\cite{chen2026mbesddpm}. Optical cloud removal~\cite{zou2023diffcr} and seismic interpolation~\cite{durall2023seismic,wei2023seismicddim} are more distant.

\subsection{Hydrology}

Diffusion models with an LSTM backbone improve extreme-flood nowcasting in most CAMELS basins~\cite{ou2025drum}. A state-space backbone with joint trajectory denoising gives more coherent ensembles~\cite{wang2025hydrodiffusion}. h-Diffusion generates hourly hydrographs from daily discharge, a form of temporal downscaling. It also assimilates observations by inpainting without retraining~\cite{yang2026hdiffusion}. A transformer-based imputer fills and forecasts sub-daily water-quantity and quality series~\cite{bhavsar2026hydrots}. We found no diffusion model for GRACE TWS among the hydrology papers we read.

\subsection{Climate emulators}

DiffESM generates daily fields consistent with monthly-mean maps; the consistency is learned rather than enforced~\cite{bassetti2024diffesm}. Spherical DYffusion emulates a global climate model probabilistically over a century~\cite{ruhlingcachay2024sphericaldyffusion}. cBottle generates global atmospheric states conditioned on monthly sea-surface temperature and reconstructs missing channels~\cite{brenowitz2025cbottle}. A global TWS generator conditioned on climate drivers would be the hydrological counterpart.

\subsection{What the intercomparisons show}

Deterministic networks regress to the mean, damp high wavenumbers and underestimate extremes. GANs and diffusion models restore them. Diffusion and flow matching are better calibrated than GANs. All but two models underestimate the climate-change signal in extremes, so visual realism does not guarantee correct trends~\cite{rampal2025intercomparison}. In a community benchmark, generative models lead clearly for precipitation, but the advantage narrows for temperature~\cite{rampal2026cordexml}. In a cross-region comparison, diffusion gains were small relative to their cost, and variability across training seeds exceeded variability across samples~\cite{singh2026intercomparison}. Earlier GAN studies showed that stochastic super-resolution can produce close to the right variability~\cite{leinonen2021gan,harris2022gan,price2022precipgen}. TWS is a smooth field, closer to temperature than to precipitation. The evidence therefore supports diffusion for GRACE mainly for coherent, calibratable ensembles, for physically plausible fine structure and for robustness to domain shift, not for large reductions in point error.

\section{Synthesis: From GRACE Tasks to Transferable Techniques}
\label{sec:synthesis}

\subsection{What distinguishes the GRACE problem}

GRACE enhancement differs from image super-resolution in four ways, each documented above.
\begin{enumerate}
  \item \emph{No fine-scale truth.} Below about 300\,km every product is a GRACE-constrained redistribution of mass according to a prior~\cite{xiong2026sysu,goracci2026nggm}. Validation there is indirect.
  \item \emph{A non-trivial observation operator.} The measurement is a truncated, filtered or regularised functional of the mass field, with correlated, anisotropic errors~\cite{kusche2007approximate,loomis2019regularization,boergens2020covariance}. It is not a block average.
  \item \emph{Very few samples.} About 250 monthly fields exist. Conditional deep models trained on so few samples risk memorisation, and the largest network can underperform~\cite{he2021deep}.
  \item \emph{Trends are the science.} Groundwater depletion and ice loss are exactly the signals that hydrological priors and climate-driven reconstructions miss~\cite{springer2026review,humphrey2019gracerec,zhang2026groundwater}.
\end{enumerate}

\subsection{Task-to-technique map}

Table~\ref{tab:map} pairs each GRACE task with the diffusion techniques that address it and the closest evidence from the works reviewed.

\begin{table*}[!t]
\centering
\caption{GRACE tasks, their specific difficulty, candidate diffusion techniques and the closest published evidence.}
\label{tab:map}
\footnotesize
\begin{tabularx}{\textwidth}{@{}l L L L@{}}
\toprule
GRACE task & Specific difficulty & Candidate techniques & Closest evidence \\
\midrule
Missing months, inter-mission gap & temporal inpainting; error regime changes at the gap & mask-conditioned and spatio-temporal diffusion; trajectory-level score models & \cite{lugmayr2022repaint,tashiro2021csdi,rozet2023sda,kishikawa2025precipinpaint,barth2024dindiff} \\
Reconstruction before 2002 & conditioning on climate forcing and sparse SLR/Swarm; trends & prior on long model simulations + likelihood guidance with observation errors & \cite{qian2026gencrest,song2025reconmost,uz2024highres,richter2021swarm_reconstruction} \\
Spatial downscaling & no truth; prior dependence; conservation & residual diffusion; guided priors with the product's operator; hard constraints at inference & \cite{mardani2025corrdiff,lopezgomez2025r2d2,tie2026zssd,hess2025scaleadaptive,bochow2026smb,harder2023hard} \\
Destriping and denoising & correlated noise of known structure; linear operator & SVD-domain and pseudoinverse posterior samplers; restoration bridges & \cite{kawar2022ddrm,song2023pigdm,chung2023dps,liu2023i2sb} \\
Temporal disaggregation & monthly to sub-monthly; consistency with the monthly mean & space--time diffusion conditioned on means; temporal downscaling of hydrographs & \cite{bassetti2024diffesm,yang2026hdiffusion} \\
Fusion with wells and GNSS & sparse point data of a different quantity & score-based assimilation with several observation operators & \cite{manshausen2025genda,chakraborty2025multimodal,li2026sage,pakrashi2026highplains} \\
Ensembles and UQ & per-pixel spread today; calibration untested & posterior sampling; ensemble amplification; multi-seed training & \cite{li2024seeds,price2025gencast,singh2026intercomparison,gou2026uncertainty} \\
\bottomrule
\end{tabularx}
\end{table*}

\subsection{Open problems}

\begin{enumerate}
  \item \emph{Operator-aware generation.} No learned GRACE method we read represents the true observation operator or its error covariance. Posterior samplers built for noisy linear operators~\cite{kawar2022ddrm,song2023pigdm,chung2023dps,tie2026zssd} offer a route, but none has been tested with spherical-harmonic truncation, decorrelation filters or a mascon resolution operator.
  \item \emph{Hard conservation for non-block operators.} Hard constraints are exact and often more accurate~\cite{harder2023hard}, and can be imposed at inference~\cite{bochow2026smb,christopher2024projected,cheng2025eci}. They have been built for block averages, not for the GRACE operator. They must also be tied to observed mass, not to a model's~\cite{hess2022pcgan}.
  \item \emph{Calibration.} No GRACE study we read reports CRPS, spread--skill ratios or rank histograms. Geoscience diffusion ensembles are often under-dispersive~\cite{mardani2025corrdiff,lopezgomez2025r2d2}, and seed variability can exceed sampling variability~\cite{singh2026intercomparison}. Formal GRACE covariance and error models~\cite{wahr2006accuracy,boergens2020covariance,loomis2019regularization,vandijk2014reanalysis} provide a reference.
  \item \emph{Small data.} Options include pretraining on long hydrological-model or closed-loop simulations~\cite{dobslaw2015esa_esm}, training from incomplete data~\cite{barth2024dindiff}, pretrained encoders~\cite{ribeiro2026prithvi}, and holding out data from the base model in residual designs~\cite{couairon2024arches}. Each brings a domain shift that must be measured~\cite{tie2026zssd}.
  \item \emph{Attribution by scale.} No study separates skill above and below the effective resolution, or the share of fine-scale variance coming from GRACE versus the prior. Spectral evaluation~\cite{lippe2023pderefiner,mardani2025corrdiff} and closed-loop truth~\cite{goracci2026nggm} make this possible.
  \item \emph{Human signals and trends.} Generative priors trained on hydrological models inherit their missing groundwater abstraction. Change signals in extremes are underestimated by most generative downscalers~\cite{rampal2025intercomparison}. Trends need explicit evaluation.
  \item \emph{Independent validation.} GNSS loading~\cite{carlson2022gnss_grace,wei2025cnn_gnss_loading}, wells~\cite{pascal2022evaluating,suryawanshi2026efficacy}, ocean-bottom pressure~\cite{chambers2012obp} and ice-mass budgets~\cite{velicogna2020continuity} are available but rarely used to test learned fine scales.
  \item \emph{Leak-free benchmarks.} A common benchmark with sequential splits, closed-loop truth and probabilistic scores across tree, U-Net, GAN and diffusion baselines does not yet exist.
  \item \emph{Future missions.} Learned methods should transfer to MAGIC-type inputs, which carry more information~\cite{goracci2026nggm,pfeffer2026nggm}.
\end{enumerate}

\subsection{Related in-house work}

In a related problem we generate daily GNSS position residual series at locations without a station. A conditional diffusion model, conditioned on neighbouring stations through cross-attention with relative-offset encodings, is evaluated with proper scores, calibration, spectra and a classifier two-sample test against kriging and Gaussian-process baselines~\cite{imonet}. Its emphasis on generators judged as samples rather than point predictions, on blind and leak-free evaluation, and on strong Gaussian baselines carries over to GRACE. Its data, which are scattered stations and daily series, do not.

% Proposed methodology: intentionally omitted until it has been discussed.
\section{Conclusion}
\label{sec:conclusion}

Learned enhancement of GRACE has matured quickly: gap fillers with uncertainty, century-long reconstructions, global 0.5$^\circ$ downscaled products and near-real-time predictions. These methods share three features. Their outputs are deterministic or carry per-pixel spread. Their fine-scale content is inherited from a hydrological or covariate prior. And they relate to GRACE through block averages and soft penalties, not through its actual measurement operator. The diffusion literature in neighbouring fields has developed what is missing: residual designs for small paired datasets, posterior sampling with explicit operators, hard constraints imposed at inference, training from incomplete data, and evaluation protocols for calibrated ensembles. The expected gain for GRACE lies less in lower point error, since TWS is smooth, than in an honest account of what gravimetry observed and what the prior supplied, expressed as a calibrated distribution of fields.


\bibliographystyle{IEEEtran}
\bibliography{references}

\end{document}
